\underline{Exercise 9.2}
\begin{enumerate}
% Question 1
\item[1.] Find the stationary values of the following (check whether they are relative maxima or minima or inflection points), assuming the domain to be the set of all real numbers:
\begin{enumerate}
\item[(c)] $y = 3x^{2} + 3$
\end{enumerate}

% Question 2
\item[2.] Find the stationary values of the following (check whether they are relative maxima or minima or inflection points), assuming the domain to be the interval $[0, \infty)$:
\begin{enumerate}
\item[(a)] $y = x^{3} - 3x + 5$
\end{enumerate}

% Question 3
\item[3.] Show that the function $y = x + 1/x$ (with $x \neq 0$) has two relative extrema, one a maximum and the other a minimum. Is the ``minimum'' larger or smaller than the ``maximum''? How is this paradoxical result possible?

% Question 4
\item[4.] Let $T = \phi(x)$ be a total function (e.g., total product or total cost):
\begin{enumerate}
\item[(a)] Write out the expressions for the marginal function $M$ and the average function $A$.
\item[(b)] Show that, when $A$ reaches a relative extremum, $M$ and $A$ must have the same value.
\item[(c)] What general principle does this suggest for the drawing of a marginal curve and an average curve in the same diagram?
\item[(d)] What can you conclude about the elasticity of the total function $T$ at the point where $A$ reaches an extreme value?
\end{enumerate}

\end{enumerate}
